\section{Analyse der Schwachstelle}
\label{sec:analyse}

\subsection{Projektstruktur}

Das Open-Source-Projekt \emph{dnchub-app} (FleetOptima) ist ein mandantenfähiges
Flottenmanagement-System auf GitHub \cite{github_issue_5}. Das Python/FastAPI-Backend
ist klassisch geschichtet:

\begin{itemize}
    \item \textbf{API-Layer:} FastAPI-Router unter \texttt{app/api/v1/}
    \item \textbf{Service-Layer:} Geschäftslogik in \texttt{app/shared/services/}
    \item \textbf{Model-Layer:} SQLAlchemy-Modelle in \texttt{app/models/}
\end{itemize}

Das gesamte Mandanten-Datenmodell basiert auf einem \texttt{organization\_id}-Feld,
das in allen geschäftsrelevanten Tabellen vorhanden ist und die logische Trennung
zwischen Mandanten sicherstellen soll.

\subsection{Phase 1: Codeanalyse}

\subsubsection{Schritt 1: Analyse der Klasse \texttt{BaseService}}

Die zentrale Klasse \texttt{BaseService} befindet sich in
\texttt{backend/app/shared/services/base.py} und implementiert sieben
Standard-CRUD-Operationen, die von allen Modulen des Systems geerbt werden.
Beim Durchsehen der Klasse fällt eine konsistente Inkonsistenz auf:
die Listenmethoden (\texttt{get\_multi}, \texttt{count}) implementieren korrekt
sowohl einen Mandanten- als auch einen Soft-Delete-Filter, die
Einzelzugriffsmethoden (\texttt{get}, \texttt{get\_or\_404}) hingegen nicht.

Tabelle \ref{tab:baseservice_audit} zeigt den Sicherheitsstatus aller Methoden:

\begin{table}[H]
    \begin{center}
        \renewcommand{\arraystretch}{1.4}
        \begin{tabular}{|l|c|c|p{4.5cm}|}
            \hline
            \textbf{Methode}        & \textbf{IDOR-sicher}           & \textbf{Soft-Delete}           & \textbf{Anmerkung}                    \\
            \hline
            \texttt{get()}          & \textcolor{red}{Nein}          & \textcolor{red}{Nein}          & Nur \texttt{WHERE id = ?}             \\
            \texttt{get\_or\_404()} & \textcolor{red}{Nein}          & \textcolor{red}{Nein}          & Delegiert an \texttt{get()}           \\
            \texttt{get\_multi()}   & \textcolor{green!60!black}{Ja} & \textcolor{green!60!black}{Ja} & \texttt{hasattr}-Guards korrekt       \\
            \texttt{count()}        & \textcolor{green!60!black}{Ja} & \textcolor{green!60!black}{Ja} & \texttt{hasattr}-Guards korrekt       \\
            \texttt{create()}       & --                             & --                             & Nicht betroffen                       \\
            \texttt{update()}       & --                             & --                             & Nicht betroffen                       \\
            \texttt{delete()}       & \textcolor{red}{Nein}          & --                             & Ruft ungesichertes \texttt{get()} auf \\
            \hline
        \end{tabular}
    \end{center}
    \caption{Sicherheitsaudit der Methoden in \texttt{BaseService} (Phase 1, Schritt 1)}
    \label{tab:baseservice_audit}
\end{table}

Listing \ref{lst:base_vulnerable} zeigt den verwundbaren Code. In \texttt{get()} wird
ausschliesslich nach \texttt{id} gefiltert, weder \texttt{organization\_id} noch
\texttt{deleted\_at} werden geprüft. Da \texttt{get\_or\_404()} intern \texttt{get()}
aufruft, erbt es beide Schwachstellen gleich mit.

\begin{lstlisting}[style=custompython,
  caption={Verwundbare Methoden in \texttt{base.py} -- Zeilen 24--34},
  label={lst:base_vulnerable}]
def get(self, db: Session, id: str) -> ModelType | None:
    """Get a single record by ID."""
    # PROBLEM: kein organization_id-Filter, kein deleted_at-Filter
    result = db.execute(select(self.model).where(self.model.id == id))
    return result.scalar_one_or_none()

def get_or_404(self, db: Session, id: str) -> ModelType:
    """Get a single record by ID or raise NotFoundError."""
    # PROBLEM: delegiert an get() -- erbt beide Schwachstellen
    obj = self.get(db, id)
    if obj is None:
        raise NotFoundError(self.model.__name__, id)
    return obj
\end{lstlisting}

Zum Vergleich zeigt Listing \ref{lst:base_correct} das Referenzmuster aus
\texttt{get\_multi()}, das beide Filter korrekt per \texttt{hasattr}-Guard implementiert:

\begin{lstlisting}[style=custompython,
  caption={Korrektes Referenzmuster aus \texttt{get\_multi()} -- Zeilen 36--57},
  label={lst:base_correct}]
def get_multi(self, db, *, organization_id=None,
              include_deleted=False):
    query = select(self.model)
    # Korrekt: Mandanten-Scoping mit Guard
    if organization_id and hasattr(self.model, "organization_id"):
        query = query.where(
            self.model.organization_id == organization_id)
    # Korrekt: Soft-Delete-Filter mit Guard
    if not include_deleted and hasattr(self.model, "deleted_at"):
        query = query.where(self.model.deleted_at.is_(None))
    ...
\end{lstlisting}

Ein weiterer sicherheitsrelevanter Befund: Auch \texttt{delete()} ruft intern das
ungesicherte \texttt{get()} auf. Ein Angreifer könnte damit nicht nur Datensätze
fremder Mandanten lesen, sondern diese auch soft-löschen, und das ohne jede
\texttt{organization\_id}-Prüfung.

\subsubsection{Schritt 2: Identifikation aller Aufrufer von \texttt{get\_or\_404()}}

Eine Volltextsuche im gesamten Backend (\texttt{grep -r "get\_or\_404"}) ergibt
\textbf{83 Aufrufstellen} in 20 Quelldateien. Tabelle \ref{tab:callers_summary}
fasst die Verteilung nach Modul und Kategorie zusammen.

\begin{table}[H]
    \begin{center}
        \renewcommand{\arraystretch}{1.4}
        \begin{tabular}{|l|r|l|}
            \hline
            \textbf{Datei / Modul}                               & \textbf{Aufrufe} & \textbf{Besonderheit}                     \\
            \hline
            \texttt{shared/api/organizations.py}                 & 5                & Wurzelelement, kein \texttt{org\_id}-Feld \\
            \texttt{shared/api/users.py}                         & 2                & Admin-only, aber Org-Scoping fehlt        \\
            \texttt{shared/api/notifications.py}                 & 2                & --                                        \\
            \texttt{shared/api/documents.py}                     & 4                & --                                        \\
            \texttt{shared/api/employees.py}                     & 3                & --                                        \\
            \texttt{shared/api/cost\_centers.py}                 & 5                & 2 Services betroffen                      \\
            \hline
            \texttt{modules/fleet/api/vehicles.py}               & 5                & 2 Services betroffen                      \\
            \texttt{modules/fleet/api/maintenance.py}            & 7                & 2 Services betroffen                      \\
            \texttt{modules/fleet/api/gps.py}                    & 7                & 3 Services betroffen                      \\
            \texttt{modules/fleet/api/tickets.py}                & 4                & --                                        \\
            \texttt{modules/fleet/api/fuel.py}                   & 2                & --                                        \\
            \texttt{modules/fleet/api/fuel\_pumps.py}            & 2                & --                                        \\
            \texttt{modules/fleet/api/fuel\_pump\_deliveries.py} & 2                & --                                        \\
            \texttt{modules/fleet/services/fuel\_pump.py}        & 4                & \texttt{self.get\_or\_404()}              \\
            \texttt{modules/fleet/services/fuel*.py}             & 2                & Cross-Service-Aufruf                      \\
            \hline
            \texttt{modules/tools/api/tools.py}                  & 3                & --                                        \\
            \texttt{modules/tools/api/tool\_cases.py}            & 4                & --                                        \\
            \texttt{modules/tools/api/consumables.py}            & 3                & --                                        \\
            \texttt{modules/tools/api/tool\_*.py}                & 7                & 4 Dateien                                 \\
            \texttt{modules/tools/services/tool\_assignment.py}  & 1                & \texttt{self.get\_or\_404()}              \\
            \hline
            \texttt{modules/sat/api/*.py}                        & 9                & 5 Dateien                                 \\
            \hline
            \textbf{Total}                                       & \textbf{83}      &                                           \\
            \hline
        \end{tabular}
    \end{center}
    \caption{Verteilung der \texttt{get\_or\_404()}-Aufrufstellen nach Modul (Phase 1, Schritt 2)}
    \label{tab:callers_summary}
\end{table}

Aus der Analyse ergeben sich drei Kategorien mit unterschiedlichem Handlungsbedarf:

\begin{enumerate}
    \item \textbf{Kategorie A: Kein Fix notwendig (\texttt{Organization}-Modell)}\\
          \texttt{organizations.py} greift auf das \texttt{Organization}-Modell zu, das
          das Wurzelelement der Mandantenhierarchie ist und selbst kein
          \texttt{organization\_id}-Attribut besitzt. Endpunkte sind entweder
          Admin-only oder lesen die eigene Organisation via
          \texttt{current\_user.organization\_id}; kein IDOR-Risiko.

    \item \textbf{Kategorie B: Partieller Fix (Admin-only-Endpunkte, \texttt{users.py})}\\
          Die beiden Aufrufe in \texttt{users.py} (Zeilen 101, 113) sind durch
          \texttt{AdminUserDep} geschützt; ein regulärer Nutzer kann sie nicht
          erreichen. Dennoch sollte für Defense-in-Depth ein Org-Scoping ergänzt werden,
          um Cross-Tenant-Zugriffe auch bei künftigen Rechteänderungen auszuschliessen.

    \item \textbf{Kategorie C: Vollständiger Fix (alle übrigen 76 Aufrufstellen)}\\
          Alle anderen Endpunkte verarbeiten mandantenspezifische Ressourcen und
          übergeben dem Aufruf kein \texttt{organization\_id}-Argument. Jeder dieser
          Aufrufe muss auf die neue Signatur umgerüstet werden und
          \texttt{current\_user.organization\_id} übergeben.
\end{enumerate}

Auffällig dabei: In den Service-Klassen \texttt{fuel\_pump.py} (4x) und
\texttt{tool\_assignment.py} (1x) wird \texttt{self.get\_or\_404()} intern genutzt,
d.h. der Aufruf erfolgt nicht aus dem API-Layer, sondern aus einem anderen Service.
Hier muss \texttt{organization\_id} durch den Service-Aufrufer durchgereicht werden.


\subsubsection{Schritt 3: Modelle mit \texttt{organization\_id} und \texttt{deleted\_at}}

Für die korrekte Implementierung des Fixes muss bekannt sein, welche Modelle welche
Felder tragen, da \texttt{BaseService} generisch über alle Modelltypen operiert.
Die Analyse aller 31 SQLAlchemy-Modell-Klassen ergibt die Klassifikation in
Tabelle \ref{tab:model_matrix}.

\begin{table}[H]
    \begin{center}
        \renewcommand{\arraystretch}{1.3}
        \begin{tabular}{|l|c|c|l|}
            \hline
            \textbf{Modell-Klasse}        & \textbf{org\_id}               & \textbf{deleted\_at}           & \textbf{Kategorie} \\
            \hline
            \multicolumn{4}{|l|}{\textit{Modul: fleet}}                                                                          \\
            \hline
            Vehicle, VehicleGroup         & \textcolor{green!60!black}{Ja} & \textcolor{green!60!black}{Ja} & Fix A              \\
            FuelEntry, FuelPump           & \textcolor{green!60!black}{Ja} & \textcolor{green!60!black}{Ja} & Fix A              \\
            FuelPumpDelivery              & \textcolor{green!60!black}{Ja} & \textcolor{green!60!black}{Ja} & Fix A              \\
            MaintenanceTask / Schedule    & \textcolor{green!60!black}{Ja} & \textcolor{green!60!black}{Ja} & Fix A              \\
            Ticket, Geofence              & \textcolor{green!60!black}{Ja} & \textcolor{green!60!black}{Ja} & Fix A              \\
            Trip                          & \textcolor{green!60!black}{Ja} & \textcolor{red}{Nein}          & Fix B              \\
            GpsAlert                      & \textcolor{green!60!black}{Ja} & \textcolor{red}{Nein}          & Fix B              \\
            VehicleGroupMember            & \textcolor{red}{Nein}          & \textcolor{red}{Nein}          & Kein Fix           \\
            TripPosition                  & \textcolor{red}{Nein}          & \textcolor{red}{Nein}          & Kein Fix           \\
            \hline
            \multicolumn{4}{|l|}{\textit{Modul: tools}}                                                                          \\
            \hline
            Tool, ToolCase, ToolCategory  & \textcolor{green!60!black}{Ja} & \textcolor{green!60!black}{Ja} & Fix A              \\
            ToolLocation, Consumable      & \textcolor{green!60!black}{Ja} & \textcolor{green!60!black}{Ja} & Fix A              \\
            ToolAssignment                & \textcolor{green!60!black}{Ja} & \textcolor{red}{Nein}          & Fix B              \\
            ToolCalibration               & \textcolor{green!60!black}{Ja} & \textcolor{red}{Nein}          & Fix B              \\
            \hline
            \multicolumn{4}{|l|}{\textit{Modul: sat}}                                                                            \\
            \hline
            SatAssistance, SatMachine     & \textcolor{green!60!black}{Ja} & \textcolor{green!60!black}{Ja} & Fix A              \\
            SatCustomer, SatContact       & \textcolor{green!60!black}{Ja} & \textcolor{green!60!black}{Ja} & Fix A              \\
            SatInterventionReport         & \textcolor{green!60!black}{Ja} & \textcolor{green!60!black}{Ja} & Fix A              \\
            SatAttachment, SatServiceType & \textcolor{green!60!black}{Ja} & \textcolor{green!60!black}{Ja} & Fix A              \\
            SatSpecialization             & \textcolor{green!60!black}{Ja} & \textcolor{green!60!black}{Ja} & Fix A              \\
            SatEmployeeSpecialization     & \textcolor{red}{Nein}          & \textcolor{red}{Nein}          & Kein Fix           \\
            \hline
            \multicolumn{4}{|l|}{\textit{Modul: shared}}                                                                         \\
            \hline
            Employee, User                & \textcolor{green!60!black}{Ja} & \textcolor{green!60!black}{Ja} & Fix A              \\
            Document, CostCenter          & \textcolor{green!60!black}{Ja} & \textcolor{green!60!black}{Ja} & Fix A              \\
            CostAllocation                & \textcolor{green!60!black}{Ja} & \textcolor{red}{Nein}          & Fix B              \\
            Organization                  & \textcolor{red}{Nein}          & \textcolor{green!60!black}{Ja} & Kein Fix*          \\
            Notification                  & \textcolor{red}{Nein}          & \textcolor{red}{Nein}          & Kein Fix           \\
            NotificationPreferences       & \textcolor{red}{Nein}          & \textcolor{red}{Nein}          & Kein Fix           \\
            \hline
        \end{tabular}
    \end{center}
    \caption{Modell-Klassifikation nach Feldern (Phase 1, Schritt 3). Fix A: beide Filter; Fix B: nur org\_id-Filter; * Organization ist Root-Entity ohne org\_id.}
    \label{tab:model_matrix}
\end{table}

Aus der Tabelle ergeben sich folgende Schlussfolgerungen für die Implementierung:

\begin{enumerate}
    \item \textbf{Fix A (27 Modelle):} Beide Filter (\texttt{organization\_id} und
          \texttt{deleted\_at}) müssen in \texttt{get()} und \texttt{get\_or\_404()} aktiv sein.
          Dies betrifft den grössten Teil aller betroffenen Endpunkte.
    \item \textbf{Fix B (5 Modelle):} Nur der \texttt{organization\_id}-Filter ist
          anzuwenden. Da diese Modelle kein \texttt{deleted\_at}-Feld besitzen, entfällt
          Issue \#14 für sie; IDOR (Issue \#5) ist jedoch dennoch präsent.
    \item \textbf{Kein Fix notwendig (5 Klassen):} Join-Tabellen
          (\texttt{VehicleGroupMember}, \texttt{TripPosition},
          \texttt{SatEmployeeSpecialization}), benutzergebundene Tabellen
          (\texttt{Notification}, \texttt{NotificationPreferences}) und die Root-Entity
          \texttt{Organization} tragen keine \texttt{organization\_id} und sind
          nicht über generic \texttt{get\_or\_404()} exponiert.
\end{enumerate}

Das \texttt{hasattr}-Guard-Muster aus \texttt{get\_multi()} und \texttt{count()}
behandelt alle Fälle korrekt, weil es zur Laufzeit prüft, ob das jeweilige Modell das
gesuchte Attribut überhaupt besitzt. Dasselbe Muster, ohne Sonderbehandlung einzelner
Modelle, kann direkt in \texttt{get()} und \texttt{get\_or\_404()} übernommen werden.


\subsubsection{Schritt 4: Analyse der bestehenden Tests}

Die vollständige Test-Suite des Projekts liegt unter \texttt{backend/tests/} und ist
in drei Schichten gegliedert:

\begin{itemize}
    \item \textbf{Unit-/API-Tests} (\texttt{tests/test\_*.py}): Prüfen einzelne Endpunkte
          (z.\,B. \texttt{test\_auth.py}, \texttt{test\_vehicles.py}) gegen den
          Anwendungs-TestClient mit echten Datenbank-Transaktionen.
    \item \textbf{Integrationstests} (\texttt{tests/integration/}): Testen vollständige
          Use-Cases über mehrere Endpunkte hinweg, z.\,B. Create-Read-Update-Delete-Zyklen.
    \item \textbf{Smoke-Tests} (\texttt{tests/smoke/}): Status-Code-Checks gegen eine
          Live-ähnliche DB; bekannte Bugs sind mit \texttt{@pytest.mark.xfail} markiert.
\end{itemize}

\paragraph{Test-Fixtures (conftest.py).}
Alle Tests verwenden eine gemeinsame Fixture-Hierarchie:
\texttt{db\_session} rollt jede Transaktion via Savepoint zurück;
\texttt{test\_organization} legt eine einzige Testorganisation an;
\texttt{test\_admin\_user} und \texttt{test\_fleet\_manager} sind User dieser Organisation.
Die Header-Fixtures \texttt{admin\_headers} / \texttt{manager\_headers} erzeugen
JWT-Bearer-Tokens via \texttt{create\_access\_token(subject, organization\_id, role)}.
\emph{Ein zweites Mandanten-Objekt existiert in der Fixture-Suite nicht}; das ist
der Ausgangspunkt für die PoC-Tests in Schritt 5.

\paragraph{Befund in integration/test\_vehicle\_crud.py.}
Die Datei enthält den Test \texttt{test\_delete\_then\_get\_still\_returns\_soft\_deleted},
der nach einem Soft-Delete explizit \texttt{status\_code == 200} erwartet und mit dem
Kommentar \emph{``GET by ID still returns 200 (soft-deleted records are not filtered
    in get\_or\_404)''} versehen ist, ein Eingeständnis des bekannten Bugs.
Dieser Test muss nach dem Patch auf \texttt{status\_code == 404} umgestellt werden.

\paragraph{Befund in smoke/test\_tools.py.}
\texttt{test\_delete\_tool\_nonexistent\_returns\_204} bestätigt ein weiteres Symptom:
\texttt{DELETE} auf eine nicht existierende ID liefert \texttt{204} statt \texttt{404},
weil \texttt{delete()} intern \texttt{get()} ohne Soft-Delete-Filter aufruft.

\paragraph{Lücke: kein Cross-Tenant-Test.}
Keiner der bestehenden Tests prüft, ob ein User die Ressourcen einer anderen
Organisation lesen kann. Alle Fixtures operieren innerhalb derselben
\texttt{test\_organization}. Diese Lücke wird in Schritt 5 durch dedizierte PoC-Tests
geschlossen.

\subsubsection{Schritt 5: Proof-of-Concept-Tests (vor dem Patch)}

Um die Schwachstellen reproduzierbar zu beweisen, wurden vier automatisierte
PoC-Tests in der Datei \texttt{tests/test\_idor\_poc.py} erstellt.
Die Tests beschreiben das \emph{korrekte, sichere} Verhalten und
\textbf{schlagen vor dem Patch fehl} (\texttt{FAILED}).
Nach dem Patch laufen sie durch (\texttt{PASSED}).
Dieses Red/Green-Muster ist der Standard im Security-Testing.

\paragraph{Zusätzliche Fixtures.}
Da alle bestehenden Fixtures nur eine einzige Organisation verwenden, wurden drei
neue Fixtures hinzugefügt:
\begin{itemize}
    \item \texttt{second\_organization}: legt eine zweite Organisation
          (\emph{Angreifer-Tenant}) direkt via SQLAlchemy an;
    \item \texttt{attacker\_user}: Admin-User in der Angreifer-Organisation;
    \item \texttt{attacker\_headers}: JWT-Bearer-Token für diesen User
          (\texttt{organization\_id} zeigt auf den Angreifer-Tenant).
\end{itemize}

\paragraph{PoC 1: IDOR (Issue \#5).}
\texttt{TestIDOR::test\_cross\_tenant\_vehicle\_read\_returns\_404}:
\begin{enumerate}
    \item Ein Fahrzeug wird der Opfer-Organisation zugeordnet (\texttt{victim\_vehicle}).
    \item Der Angreifer-User ruft \texttt{GET /api/v1/vehicles/\{id\}} auf.
    \item \textbf{Korrektes Verhalten (Test-Assertion):} \texttt{404 Not Found}.
    \item \textbf{Aktuelles Verhalten:} \texttt{200 OK} (Test \texttt{FAILED},
          Schwachstelle bewiesen).
\end{enumerate}

\texttt{TestIDOR::test\_cross\_tenant\_vehicle\_read\_does\_not\_expose\_foreign\_org\_id}
prüft ebenfalls, dass der Endpunkt \texttt{404} zurückliefert und damit keine
\texttt{organization\_id} des Opfer-Tenants preisgibt.

\paragraph{PoC 2: Soft-Delete-Bypass (Issue \#14).}
\texttt{TestSoftDeleteBypass::test\_soft\_deleted\_vehicle\_returns\_404}:
\begin{enumerate}
    \item Ein Fahrzeug wird via \texttt{DELETE /api/v1/vehicles/\{id\}} soft-deleted
          (Rückgabecode: \texttt{204}).
    \item \texttt{GET /api/v1/vehicles/\{id\}} durch den Admin-User.
    \item \textbf{Korrektes Verhalten:} \texttt{404 Not Found}.
    \item \textbf{Aktuelles Verhalten:} \texttt{200 OK} (Test \texttt{FAILED}).
\end{enumerate}

\texttt{TestSoftDeleteBypass::test\_soft\_deleted\_vehicle\_absent\_in\_list\_and\_by\_id}
zeigt die Inkonsistenz: \texttt{GET /api/v1/vehicles} (Liste) liefert das
soft-gelöschte Fahrzeug korrekt nicht zurück, \texttt{GET /api/v1/vehicles/\{id\}}
hingegen schon; beide Assertions müssen \texttt{404} ergeben.

\paragraph{Testergebnis vor dem Patch.}
\begin{center}
    \begin{tabular}{lll}
        \hline
        \textbf{Testname}                               & \textbf{Ergebnis (vor Patch)} & \textbf{Erwartet (nach Patch)} \\
        \hline
        \texttt{test\_cross\_tenant\_...\_returns\_404} & \texttt{FAILED}               & \texttt{PASSED}                \\
        \texttt{test\_cross\_tenant\_...\_no\_expose}   & \texttt{FAILED}               & \texttt{PASSED}                \\
        \texttt{test\_soft\_deleted\_...\_returns\_404} & \texttt{FAILED}               & \texttt{PASSED}                \\
        \texttt{test\_soft\_deleted\_...\_absent\_...}  & \texttt{FAILED}               & \texttt{PASSED}                \\
        \hline
    \end{tabular}
\end{center}

\vspace{0.5em}
\noindent Alle vier Tests scheitern erwartungsgemäss vor dem Patch
(\texttt{4 failed}) und beweisen damit die Schwachstellen reproduzierbar.
